\documentclass[11pt,a4paper]{article}

\usepackage{fontspec}
\setmainfont{DejaVu Serif}[
  BoldFont={DejaVu Serif Bold},
  ItalicFont={DejaVu Serif},
  ItalicFeatures={FakeSlant=0.18},
  BoldItalicFont={DejaVu Serif Bold},
  BoldItalicFeatures={FakeSlant=0.18}
]
\setsansfont{DejaVu Sans}
\newfontfamily\cyrillicfont{DejaVu Serif}
\newfontfamily\cyrillicfontsf{DejaVu Sans}
\newfontfamily\cyrillicfonttt{DejaVu Sans Mono}
\usepackage{polyglossia}
\setmainlanguage{english}
\usepackage[a4paper,margin=20mm]{geometry}
\usepackage{microtype}
\usepackage{setspace}
\setstretch{1.015}
\usepackage{amsmath,amssymb,amsthm,mathtools,bm}
\usepackage{booktabs,longtable,array,tabularx}
\usepackage{enumitem}
\usepackage{xcolor}
\usepackage{tikz}
\usetikzlibrary{arrows.meta,positioning,calc,fit}
\usepackage{xurl}
\usepackage{seqsplit}
\usepackage[unicode,hidelinks]{hyperref}
\usepackage[nameinlink,noabbrev]{cleveref}
\emergencystretch=3em
\allowdisplaybreaks

\hypersetup{
  pdftitle={Geometry of Observation: Frozen Typed Corpus and Reproducibility Master},
  pdfauthor={Stassis Stashkevichyus},
  pdfsubject={P12 corpus freeze of eighteen GO Core reference modules},
  pdfkeywords={geometry of observation, typed observation maps, invariance, covariance, identifiability, reproducibility, corpus freeze}
}

\definecolor{obsblue}{RGB}{37,83,138}
\definecolor{obsred}{RGB}{150,45,45}
\definecolor{obsgreen}{RGB}{30,115,78}
\definecolor{obsviolet}{RGB}{94,64,145}
\definecolor{obsgray}{RGB}{74,78,84}
\definecolor{lightblue}{RGB}{239,246,253}
\definecolor{lightred}{RGB}{253,242,242}
\definecolor{lightgreen}{RGB}{239,249,244}
\definecolor{lightviolet}{RGB}{246,242,252}

\theoremstyle{definition}
\newtheorem{definition}{Definition}[section]
\newtheorem{protocol}[definition]{Protocol}
\theoremstyle{plain}
\newtheorem{proposition}[definition]{Proposition}
\theoremstyle{remark}
\newtheorem{boundary}[definition]{Boundary}
\newtheorem{remark}[definition]{Remark}

\newcommand{\R}{\mathbb R}
\newcommand{\Prob}{\mathsf P}
\newcommand{\Hash}{\operatorname{SHA256}}
\newcommand{\id}{\operatorname{id}}
\newcommand{\rank}{\operatorname{rank}}
\newcommand{\qual}[2]{\texttt{#1::#2}}
\newcommand{\norm}[1]{\left\lVert #1\right\rVert}
\newcommand{\abs}[1]{\left\lvert #1\right\rvert}
\newcommand{\Xcal}{\mathcal X}
\newcommand{\Ycal}{\mathcal Y}
\newcommand{\Ical}{\mathcal I}
\newcommand{\Qcal}{\mathcal Q}
\newcommand{\Scal}{\mathcal S}
\newcommand{\Acal}{\mathcal A}
\newcommand{\Fcal}{\mathfrak F}
\newcolumntype{L}[1]{>{\raggedright\arraybackslash}p{#1}}
\newcolumntype{Y}{>{\raggedright\arraybackslash}X}

\newcommand{\ModuleCount}{18}
\newcommand{\ComponentPageCount}{158}
\newcommand{\TypedExpressionCount}{347}
\newcommand{\ClaimCount}{136}
\newcommand{\StrongClaimCount}{103}
\newcommand{\ModelClaimCount}{14}
\newcommand{\HypothesisClaimCount}{11}
\newcommand{\DiagnosticClaimCount}{7}
\newcommand{\EmpiricalClaimCount}{1}
\newcommand{\DependencyNodeCount}{29}
\newcommand{\DependencyEdgeCount}{36}
\newcommand{\ShortIdentifierOverlapCount}{7}
\newcommand{\FrontMatterPageCount}{8}
\newcommand{\MasterPageCount}{166}
\newcommand{\BaselineLedgerHash}{\texttt{\seqsplit{3f50f7941f8549c4b9de626fdb107431c07238582cdf8ebce90203500189f68e}}}


\title{\bfseries Geometry of Observation\\
\Large Frozen Typed Corpus and Reproducibility Master\\[0.35em]
\large Eighteen reference modules, one scoped contract, and an immutable release map}
\author{Stassis Stashkevichyus\\
\small Independent Research Program\\
\small ORCID: \href{https://orcid.org/0009-0000-2294-705X}{0009-0000-2294-705X}}
\date{P12 corpus freeze v1.3.0 --- 28 July 2026}

\begin{document}
\maketitle

\begin{center}
\fcolorbox{obsblue}{lightblue}{%
\parbox{0.92\linewidth}{\small
\textbf{Release status.}
This master freezes \ModuleCount\ normative modules with
\ComponentPageCount\ component pages, \TypedExpressionCount\ typed
expressions, and zero GO Core findings.  It is a provenance and
reproducibility object.  Concatenation does not turn the component
results into one theorem, does not establish physical equivalence
between application domains, and does not claim a new theory of
general relativity, information theory, mechanics, quantum chemistry,
accelerator physics, or astrodynamics.}}
\end{center}

\begin{abstract}
This release freezes the strict Geometry of Observation corpus after
the P0--P11 migration sequence.  The common mathematical discipline is
not a universal physical dynamics but a typed separation of hidden
state, invertible representation change, reduction, observation
channel, sampling protocol, estimator, and declared invariant or
recoverable target.  The frozen object contains eighteen independently
scoped reference modules, the GO Core v0.2 type and dimension contract,
a machine-readable corpus ledger, an acyclic dependency ledger,
canonical release metadata, exact component hashes, and executable
validation.  Every component PDF is embedded without textual revision;
its original hash remains the authority for component identity.
Short identifiers are document-scoped and receive fully qualified
release names.  Strong, model, hypothesis, diagnostic, and empirical
claims remain separated.  The master therefore records a clean corpus
state while preserving the scientific boundaries of each module.
\end{abstract}

\begin{sloppypar}
\raggedright
\noindent\textbf{Аннотация.}
Релиз фиксирует строгий корпус геометрии наблюдения после миграций
P0--P11. Общей структурой является не единая физическая динамика, а
типизированное разделение скрытого состояния, обратимой смены
представления, редукции, канала наблюдения, дискретизации, оценивателя и
явно объявленного инварианта или восстанавливаемой величины. В master
включены восемнадцать модулей без переписывания их содержимого,
машинные реестры, зависимости, контрольные суммы и воспроизводимые
проверки. Совпадение структуры между областями не трактуется как
доказательство физической эквивалентности.
\end{sloppypar}

\tableofcontents
\clearpage

\section{The frozen release object}

\begin{definition}[Corpus freeze]
Let
\begin{equation}
\Fcal
=
\bigl(C,L,O,D,S,H,M\bigr)
\end{equation}
denote, respectively, the core contract, corpus ledger, canonical
module order, dependency ledger, component source map, component hash
map, and release metadata.  A release is frozen when all seven entries
are versioned and every file named by \(S\) has a verified digest in
\(H\).  Two freezes are identical precisely when their canonical
serializations and every referenced byte stream agree.
\end{definition}

The baseline ledger is
\texttt{go-corpus-ledgers@1.2.0}, with digest
\BaselineLedgerHash.  P12 does not mutate that ledger.
Instead it adds a release-level ledger whose order and dependency
semantics are explicit.  A change to a component byte, ordering rule,
dependency edge, author identity, licence, or release metadata requires
a new release version.  A filename change alone is not sufficient to
create a new scientific object, while a changed byte stream cannot
retain the old freeze identity.

\begin{proposition}[Hash-addressed component identity]
Assume collision resistance of SHA-256 for the finite release set.
If two release records have the same ordered component paths, byte
counts, and SHA-256 digests, then the probability of an undetected
component substitution is negligible under the stated cryptographic
assumption.
\end{proposition}

\begin{boundary}
This is an integrity statement, not a mathematical proof that the
scientific claims are true.  Hash agreement proves preservation of the
audited object; it does not replace peer review.
\end{boundary}

\section{Common typed architecture}

The minimal release-wide observation pattern is
\begin{equation}\label{eq:master-chain}
\Xcal
\xrightarrow{\ \Phi_t\ }
\Xcal_t
\xrightarrow{\ F_g\ }
\Xcal_t'
\xrightarrow{\ R_\lambda\ }
\Xcal_\lambda
\xrightarrow{\ \Qcal_\lambda\ }
\Ycal_\lambda
\xrightarrow{\ \Scal_{\Delta t,T}\ }
\Ycal_\lambda^M
\xrightarrow{\ \Acal\ }
\Ical .
\end{equation}
Here \(F_g\) is an invertible frame, gauge, coordinate, or basis
change.  The reduction \(R_\lambda\), observation \(\Qcal_\lambda\),
and sampler \(\Scal_{\Delta t,T}\) may lose information.  The estimator
\(\Acal\) is a decision rule, not a physical transformation of the
hidden system.  A module may omit stages only by declaring the
corresponding identity or noiseless map.

For a declared group \(G\), protocol \(\lambda\), and representation
\(\rho\), the release distinguishes
\begin{align}
I(g\!\cdot\!x;\lambda)&=I(x;\lambda)
&&\text{(invariance)},\\
\Qcal_\lambda(g\!\cdot\!x)&=\rho(g)\Qcal_\lambda(x)
&&\text{(covariance)},\\
T(x)&=\overline T\!\left(\Qcal_\lambda(x)\right)
&&\text{(deterministic recoverability)}.
\end{align}
In a stochastic finite-alphabet setting, exact recoverability of a
target \(T(X)\) from \(Y\) is equivalently stated as
\begin{equation}
H\!\left(T(X)\mid Y\right)=0
\end{equation}
under the declared probability law.  None of these statements is
well-typed without its transformation group, domain, protocol, and
unit context.

\begin{boundary}[No cross-domain identity claim]
The recurrence of \cref{eq:master-chain} in geometry, mechanics,
spectroscopy, molecular inference, or satellite networks records a
shared inference architecture.  It does not identify their state
spaces, Hamiltonians, field equations, probability laws, or empirical
content.
\end{boundary}

\section{Corpus layers and canonical order}

The master order is semantic rather than chronological.  It preserves
every component identity while placing interfaces before the
laboratories that use them.

\begin{center}
\begin{tikzpicture}[
  node distance=7mm and 8mm,
  box/.style={draw,rounded corners,minimum width=34mm,minimum height=9mm,
    align=center,font=\small,fill=lightblue},
  app/.style={box,fill=lightgreen},
  spec/.style={box,fill=lightviolet},
  arr/.style={-{Latex[length=2mm]},thick,obsgray}
]
\node[box] (foundation) {Foundations\\typed maps and losses};
\node[box,below=of foundation] (metric) {Metric and scale\\quotients and resolution};
\node[app,below left=of metric] (mechanics) {Classical systems\\frames and contact};
\node[spec,below right=of metric] (spectral) {Spectral inference\\molecular channels};
\node[app,below=15mm of metric] (relativistic) {Relativistic and temporal systems};
\draw[arr] (foundation) -- (metric);
\draw[arr] (metric) -- (mechanics);
\draw[arr] (metric) -- (spectral);
\draw[arr] (foundation.east) to[bend left=22] (relativistic.east);
\draw[arr] (mechanics) -- (relativistic);
\end{tikzpicture}
\end{center}

\begin{longtable}{@{}r L{23mm} L{67mm} c r r r@{}}
\caption{Canonical semantic order of the frozen corpus. Counts are module-local.}\label{tab:module-catalog}\\
\toprule
No. & Layer & Module & v & pp. & Expr. & Claims\\
\midrule
\endfirsthead
\toprule
No. & Layer & Module & v & pp. & Expr. & Claims\\
\midrule
\endhead
\bottomrule
\endfoot
1 & Foundation & Functional Interface GO-QS & 0.2.0 & 12 & 3 & 2\\
2 & Foundation & Tensorial Observation Geometry in GR & 0.3.0 & 11 & 2 & 3\\
3 & Information & Information-Theoretic Observation Geometry & 0.2.0 & 9 & 9 & 9\\
4 & Metric & Metric Entropy and Observational Entropy Defect & 0.2.0 & 10 & 13 & 12\\
5 & Metric & Distance and Scale Interface under Observation Maps & 0.2.0 & 9 & 12 & 8\\
6 & Scale & Planck-to-Cosmos Observation Rulers & 1.1.0 & 12 & 19 & 10\\
7 & Scale & Mandelbrot Rulers as Observation-Scale Geometry & 1.1.0 & 11 & 12 & 9\\
8 & Discrete geometry & Regular Polyhedra under Typed Observation Filters & 1.1.0 & 7 & 42 & 10\\
9 & Mechanics & Frames, Forces, Constraints, and Dissipation under Observation Maps & 0.1.0 & 7 & 10 & 4\\
10 & Mechanics & Foucault Networks on Celestial Bodies & 1.1.0 & 6 & 7 & 4\\
11 & Dynamics & Billiards as a Typed Geometry of Observation Laboratory & 1.1.0 & 9 & 24 & 13\\
12 & Contact mechanics & Bobsleigh Contact Geometry as a Typed Observation System & 1.1.0 & 6 & 11 & 4\\
13 & Mechanics & Roller-Coaster Geometry as a Typed Observation Laboratory & 1.1.0 & 6 & 11 & 4\\
14 & Contact mechanics & Gear Contact Geometry as a Typed Finite Observation System & 1.1.0 & 9 & 20 & 9\\
15 & Spectral & Conical Intersections as Typed Spectral Singularities under Observation Maps & 1.1.0 & 8 & 36 & 11\\
16 & Molecular inference & Quantum Chemistry as a Typed Inference Stack under Observation Maps & 1.1.0 & 9 & 47 & 11\\
17 & Relativistic & Relativistic Beam Paths as Observation Geometry & 1.3.0 & 9 & 16 & 4\\
18 & Temporal networks & Satellite Networks under Typed Frames and Temporal Observation Channels & 1.2.0 & 8 & 53 & 9\\
\end{longtable}


\section{Inter-module dependency audit}

Dependencies are typed.  A \emph{core-contract} edge is normative: the
target ledger must validate against GO Core v0.2.  An
\emph{extension-contract} edge binds a module to an auxiliary
contract such as the SI--HEP passport or the mechanics interface.  A
\emph{documented-interface} edge records an explicit source-level
reference between modules.  A \emph{contextual} relation is excluded
from the normative graph and cannot be used as a proof dependency.

\begin{longtable}{@{}L{51mm} L{67mm} L{31mm}@{}}
\caption{Explicit non-global dependency edges.  The universal GO Core edge to all modules is omitted from the rows.}\label{tab:dependency-edges}\\
\toprule
From & To & Edge kind\\
\midrule
\endfirsthead
\toprule
From & To & Edge kind\\
\midrule
\endhead
\bottomrule
\endfoot
go-distance-scale-contract@0.3.0 & Distance and Scale Interface under Observation Maps & extension contract\\
go-distance-scale-contract@0.3.0 & Mandelbrot Rulers as Observation-Scale Geometry & extension contract\\
go-planck-cosmos-scale-contract@0.4.0 & Planck-to-Cosmos Observation Rulers & extension contract\\
go-si-hep-quantity-passport@0.5.0 & Relativistic Beam Paths as Observation Geometry & extension contract\\
go-frame-force-dissipation-contract@0.6.0 & Frames, Forces, Constraints, and Dissipation unde... & extension contract\\
go-frame-force-dissipation-contract@0.6.0 & Foucault Networks on Celestial Bodies & extension contract\\
go-frame-force-dissipation-contract@0.6.0 & Bobsleigh Contact Geometry as a Typed Observation... & extension contract\\
go-frame-force-dissipation-contract@0.6.0 & Roller-Coaster Geometry as a Typed Observation La... & extension contract\\
go-frame-force-dissipation-contract@0.6.0 & Gear Contact Geometry as a Typed Finite Observati... & extension contract\\
go-gear-contact-contract@0.7.0 & Gear Contact Geometry as a Typed Finite Observati... & extension contract\\
go-billiards-observation-contract@0.8.0 & Billiards as a Typed Geometry of Observation Labo... & extension contract\\
go-conical-intersections-observation-contract@0.9.0 & Conical Intersections as Typed Spectral Singulari... & extension contract\\
go-quantum-chemistry-observation-contract@1.0.0 & Quantum Chemistry as a Typed Inference Stack unde... & extension contract\\
go-regular-polyhedra-observation-contract@1.1.0 & Regular Polyhedra under Typed Observation Filters & extension contract\\
go-satellite-networks-observation-contract@1.2.0 & Satellite Networks under Typed Frames and Tempora... & extension contract\\
Information-Theoretic Observation Geometry & Metric Entropy and Observational Entropy Defect & documented interface\\
Metric Entropy and Observational Entropy Defect & Distance and Scale Interface under Observation Maps & documented interface\\
Distance and Scale Interface under Observation Maps & Planck-to-Cosmos Observation Rulers & documented interface\\
\end{longtable}


The normative and documented-interface graph has
\DependencyNodeCount\ nodes and \DependencyEdgeCount\ edges.  Its
topological sort is unique only up to independent branches; acyclicity
is the relevant release condition.  Each edge endpoint resolves to a
frozen module or a versioned contract, and each documented evidence
fragment is checked against the cited source.

\section{Namespaces and symbol discipline}

Document-local identifiers are not silently promoted to global names.
For document \(d\) and local identifier \(a\), the master name is
\begin{equation}
\operatorname{qid}(d,a)=\texttt{d::a}.
\end{equation}
This operation is injective on ordered pairs.  The audit found
\ShortIdentifierOverlapCount\ repeated short identifiers across
expression and claim registers; all become distinct after
qualification.  These overlaps are reuse of local vocabulary, not
ledger collisions.  Map identifiers and quantity identifiers are
already globally unique in the frozen ledger.

Symbols remain scoped records, not bare glyphs.  Thus \(G\) may denote
a symmetry group in one module and the gravitational constant in
another only when the scope and quantity identifier are carried with
the expression.  The master does not create a universal symbol table
by typographic coincidence.

\section{Claim stratification and scientific boundary}

The corpus claim register contains \ClaimCount\ entries:
\StrongClaimCount\ theorem/proposition/corollary records,
\ModelClaimCount\ models, \HypothesisClaimCount\ hypotheses,
\DiagnosticClaimCount\ diagnostics, and \EmpiricalClaimCount\ empirical
record.  A strong label is accepted by the release only when the local
ledger supplies the hypotheses and the source states the result in its
declared scope.  P12 does not strengthen a component claim.

\begin{center}
\begin{tabularx}{0.96\linewidth}{@{}L{49mm}Y@{}}
\toprule
Layer & Release interpretation\\
\midrule
Definition & Fixes notation or an object; no empirical claim.\\
Theorem, proposition, corollary & Mathematical statement under the local hypotheses.\\
Model & Declared idealization or computational benchmark.\\
Hypothesis & Unproved assumption or proposed condition.\\
Diagnostic & Protocol-dependent summary; not an invariant unless separately proved.\\
Empirical & Externally measured input with provenance and uncertainty requirements.\\
\bottomrule
\end{tabularx}
\end{center}

\begin{boundary}[Master claim firewall]
The release does not claim a theory of everything, a derivation of
Standard Model parameters, a new electronic-structure method, a new
ergodicity theorem, a new classification of regular polyhedra, or an
operational orbit-determination system.  Its defensible contribution
is the typed audit infrastructure and the module-specific corrections
recorded in the appended documents.
\end{boundary}

\section{Metadata normalization and provenance}

The canonical release author is \textbf{Stassis Stashkevichyus},
identified by ORCID
\href{https://orcid.org/0009-0000-2294-705X}
{0009-0000-2294-705X}.  Earlier component PDFs contain the historical
display strings \emph{Stassis Research Program} or
\emph{Stas, Independent Research Program}.  Their bytes are preserved;
the container metadata supplies one normalized creator record without
pretending that the old PDF metadata was rewritten.

No DOI is fabricated in this freeze.  The release metadata is prepared
for one corpus-level archival record; component modules are not assigned
separate placeholder identifiers.  If a DOI is minted, the immutable
release metadata must be updated in a successor version or by an
archive-controlled identifier field that does not alter the deposited
files.

The root citation file follows Citation File Format 1.2.0.  The Zenodo
metadata file follows the documented GitHub-release override structure.
The release is marked CC BY-NC-ND 4.0 as directed by the author.  This
licence permits non-commercial redistribution of unadapted material
with attribution; it is not an open-source software licence and does
not permit distribution of modified versions.

\section{Validation and reproducibility}

The P12 gate requires all of the following:
\begin{enumerate}[leftmargin=*,itemsep=2pt]
  \item exact resolution of all \ModuleCount\ PDF and source paths;
  \item agreement of component page counts, byte counts, and SHA-256
    digests with the freeze ledger;
  \item a clean GO Core replay:
    \(18\,\mathrm{PASS}/0\,\mathrm{FAIL}/0\,\mathrm{BLOCKED}\);
  \item uniqueness of document IDs and fully qualified local IDs;
  \item an acyclic, endpoint-complete dependency graph with verified
    source evidence;
  \item successful execution of the P1--P11 release validators;
  \item A4 geometry, no encryption or form fields, embedded fonts, and
    extractable text for every component and the master;
  \item page-content preservation between each source PDF and its
    corresponding master page interval;
  \item valid citation, archival, licence, and checksum records;
  \item a deterministic ZIP archive with fixed timestamps, sorted
    members, and byte-identical rebuild.
\end{enumerate}

The master contains \FrontMatterPageCount\ release-front-matter pages
followed by the \ComponentPageCount\ unchanged component pages, for a
total of \MasterPageCount\ pages.  PDF bookmarks and the external page
map identify every component boundary.  The component PDFs, not the
master page numbering, remain the primary unit of scientific citation.

\section{How to cite and review this object}

Scientific review should proceed module by module.  A reviewer should
first inspect the claim register and protocol fields, then the local
proofs or counterexamples, and finally the numerical and release
checks.  A criticism of one application module does not automatically
invalidate the typed core; conversely, a clean core lint does not
validate an application theorem whose hypotheses are false.

Until an archival DOI exists, cite the title, author, ORCID, version,
date, and the SHA-256 digest recorded in
\texttt{SHA256SUMS\_v1\_3.txt}:
\begin{quote}\small
Stassis Stashkevichyus,
\emph{Geometry of Observation: Frozen Typed Corpus and Reproducibility
Master}, P12 corpus freeze v1.3.0, 28 July 2026,
ORCID 0009-0000-2294-705X.
\end{quote}

\section{Release conclusion}

P12 closes the migration sequence at the corpus level.  The defensible
result is an immutable, typed, hash-addressed collection of eighteen
reference modules with explicit claim boundaries and reproducible
validation.  The release is suitable for archival provenance and
external module-by-module review.  It is not, by itself, evidence that
the umbrella programme has become a single new physical theory.

\end{document}
